% =============================================================================
% SOB STRATEGIEMODELL - OFFERING (OFFERTE)
% =============================================================================
% FehrAdvice & Partners AG
% Unterschriftsreife Offerte für Besprechung mit Armin Weber
% Version 3.0: Angereichert mit diagnostischen Highlights aus 11-Dokument-Analyse
% Compile with: xelatex SOB002_strategiemodell_offering.tex
% =============================================================================

\documentclass{fehradvice-report}

% =============================================================================
% DOCUMENT METADATA
% =============================================================================

\title{Offerte: SOB Strategiemodell}
\subtitle{Projekt SOB002 -- Evidenzbasierte Strategieentwicklung}
\author{FehrAdvice \& Partners AG}
\date{10. Februar 2026}

\client{Schweizerische Südostbahn AG}
\project{SOB002}
\version{3.0}
\classification{Vertraulich}

% =============================================================================
% DOCUMENT
% =============================================================================

\begin{document}

% Title page
\maketitle

\vspace{1cm}

% =============================================================================
\section*{Offerte Nr. SOB002-2026-01}
% =============================================================================

\begin{center}
\begin{tabular}{p{4cm}p{9cm}}
\textbf{Auftraggeber:} & Schweizerische Südostbahn AG \\
                       & Bahnhofplatz 1a \\
                       & 9001 St. Gallen \\
                       & z.Hd. Armin Weber, CEO \\[0.3cm]
\textbf{Auftragnehmer:} & FehrAdvice \& Partners AG \\
                        & Klausstrasse 4 \\
                        & 8008 Zürich \\[0.3cm]
\textbf{Projekt:} & SOB Strategiemodell (SOB002) \\
\textbf{Referenz:} & Debriefing-Gespräch vom 28. Januar 2026 \\[0.3cm]
\textbf{Offerten-Datum:} & 10. Februar 2026 \\
\end{tabular}
\end{center}

\newpage
% =============================================================================
\section{Ausgangslage und Herausforderung}
% =============================================================================

\subsection{Erfolgreicher Abschluss SOB001}

Mit dem Projekt SOB001 (Oktober 2025 -- Januar 2026) haben wir gemeinsam wichtige Grundlagen geschaffen:

\begin{itemize}
    \item \textbf{Neues Mission Statement:} <<Einsteigen -- Gemeinsam Bewegen -- Wege bahnen>> (Score 4.26/5.0)
    \item \textbf{Kontextvektor:} 32 verhaltensrelevante Faktoren für die SOB dokumentiert
    \item \textbf{Verhaltensparameter:} Segmentspezifische Parameter kalibriert
\end{itemize}

\subsection{Vertiefte Analyse: Was uns 11 Dokumente verraten}

In Vorbereitung auf SOB002 haben wir die 11 zentralen Strategiedokumente der SOB (2012--2025) umfassend analysiert. Die Befunde zeigen klare Handlungsfelder:

\begin{center}
\begin{tabular}{p{4cm}p{2.5cm}p{6.5cm}}
\toprule
\rowcolor{fadarkblue}
\fahead{Befund} & \fahead{Kennzahl} & \fahead{Was das bedeutet} \\
\midrule
\rowcolor{fawhite}
\textbf{Strategie-Lücke} & Pyramide 31\% & Vision und Werte sind da, aber strategische Ziele, Massnahmen und KPIs fehlen (Ebenen 6--8) \\
\rowcolor{falightgray}
\textbf{Integrations-Lücke} & $\gamma = 0.27$ & Die GBs arbeiten gut für sich, aber die Zusammenarbeit ist deutlich unter dem Potenzial (Ziel: 0.52) \\
\rowcolor{fawhite}
\textbf{Umsetzungs-Lücke} & Say-Do Gap 9 & Werte sind bekannt (Score 78), aber nicht gelebt (Score 69) \\
\rowcolor{falightgray}
\textbf{Steuerungs-Lücke} & 0 Leading KPIs & Die GL hat kein operatives Dashboard zur Strategiesteuerung \\
\bottomrule
\end{tabular}
\end{center}

\begin{fainfobox}[Die zentrale Erkenntnis]
Durch bessere Vernetzung der Geschäftsbereiche ($\gamma$ von 0.27 auf 0.52) kann die SOB ihre Gesamtleistung um rund \textbf{25\%} steigern -- ohne dass ein einzelner Bereich mehr leisten muss. Das ist die grösste ungenutzte Chance.
\end{fainfobox}

\subsection{Zeitdruck: Konzessionserneuerung 2029}

Die Konzessionserneuerung 2029 erfordert, dass die SOB bis 2028 ein nachweisbar funktionierendes Strategiemodell vorweisen kann. Das Zeitfenster für SOB002 ist daher jetzt.

\newpage
% =============================================================================
\section{Zielsetzung}
% =============================================================================

\begin{fainfobox}[Ziel des Projekts]
Entwicklung eines integrierten SOB-Strategiemodells in BEATRIX (EBF), das die 47 identifizierten Lücken systematisch adressiert und die Strategieumsetzung messbar macht.
\end{fainfobox}

\subsection{Messbare Kernziele}

\begin{center}
\begin{tabular}{p{0.5cm}p{4cm}p{3cm}p{3cm}p{2.5cm}}
\toprule
\rowcolor{fadarkblue}
\fahead{\#} & \fahead{Ziel} & \fahead{IST} & \fahead{SOLL} & \fahead{Messung} \\
\midrule
\rowcolor{fawhite}
Z1 & Strategie-Pyramide vervollständigen & 31\% & 100\% & Pyramide-Audit \\
\rowcolor{falightgray}
Z2 & Bereichs-Komplementarität erhöhen & $\gamma = 0.27$ & $\gamma = 0.35$ & $\gamma$-Messung \\
\rowcolor{fawhite}
Z3 & Steuerungsfähigkeit herstellen & 0 Leading KPIs & $\geq$5 & KPI-Dashboard \\
\bottomrule
\end{tabular}
\end{center}

\subsection{Erwartete Ergebnisse}

\begin{center}
\begin{tabular}{p{4cm}p{9cm}}
\toprule
\rowcolor{fadarkblue}
\fahead{Für wen} & \fahead{Nutzen} \\
\midrule
\rowcolor{fawhite}
SOB-Führung (GL) & Evidenzbasiertes Strategiemodell mit messbaren Dimensionen und Fortschritts-Tracking \\
\rowcolor{falightgray}
Bereiche (5 GBs) & Bereichsspezifische Factsheets mit konkreten $\gamma$-Zielen für die Zusammenarbeit \\
\rowcolor{fawhite}
Organisation & Konsistente Strategiearchitektur von der Vision bis zur operativen Umsetzung \\
\rowcolor{falightgray}
Verwaltungsrat & Präsentation des Gesamtmodells -- Grundlage für die Konzessionserneuerung 2029 \\
\bottomrule
\end{tabular}
\end{center}

\newpage
% =============================================================================
\section{Unser geplantes Vorgehen}
% =============================================================================

Das Projekt gliedert sich in drei Phasen über 12 Wochen (Februar -- April 2026):

\subsection{Phase 1: Konsolidierung (Wochen 1--4)}

\textbf{Ziel:} Strategie-Architektur vervollständigen und mit GL validieren.

\begin{fainfobox}[Vorsprung durch Voranalyse]
Folgende Analysen sind bereits erarbeitet und fliessen direkt in Phase~1 ein:
\begin{itemize}
    \item Strategie-Pyramide IST-Analyse (8 Ebenen, 31\% Vollständigkeit)
    \item Stakeholder-Mapping (14 Akteure, Multi-Prinzipal-Dynamik)
    \item Gap-Analyse (47 Gaps in 7 Dimensionen, priorisiert)
    \item Bereichs-Komplementarität (5$\times$5 $\gamma$-Matrix)
\end{itemize}
\end{fainfobox}

\begin{itemize}
    \item \textbf{Strategie-Pyramide vervollständigen:} Ebenen 6--8 (Ziele, Massnahmen, KPIs) gemeinsam mit GL definieren
    \item \textbf{Gap-Validierung mit GL:} Die 47 identifizierten Gaps priorisieren und Handlungsfelder festlegen
    \item \textbf{Bereichsleiter-Gespräche:} Individuelle Perspektiven auf Strategie und Zusammenarbeit erfassen
\end{itemize}

\textbf{Deliverable:} Vollständige Strategie-Pyramide (1 Seite) mit GL-Freigabe

\subsection{Phase 2: Modellentwicklung (Wochen 5--8)}

\textbf{Ziel:} Vollständiges Strategiemodell in BEATRIX/EBF entwickeln.

\begin{itemize}
    \item \textbf{Strategie-Dimensionen definieren:} Basierend auf den 3 Kerndefiziten und der $\gamma$-Matrix die relevanten Dimensionen ableiten
    \item \textbf{Parameter kalibrieren:} Verhaltensparameter pro Dimension -- fundiert auf MA-Umfrage (N=771), Strategiedokumenten und EBF-Literatur
    \item \textbf{$\gamma$-Interventionen modellieren:} Welche Massnahmen erhöhen die Bereichs-Komplementarität am effektivsten?
    \item \textbf{Modell in BEATRIX implementieren:} SOB-Strategiemodell aufbauen und mit Kontextvektor verknüpfen
\end{itemize}

\textbf{Deliverables:} SOB-Strategiemodell in BEATRIX, Dimensions-Katalog

\subsection{Phase 3: Implementierungsplanung (Wochen 9--12)}

\textbf{Ziel:} Strategie für die Umsetzung in den Bereichen vorbereiten.

\begin{itemize}
    \item \textbf{$\gamma$-Dashboard:} Interaktives Tool, das die Komplementarität zwischen GBs laufend sichtbar macht
    \item \textbf{Bereichsübersetzung:} Für alle 5 GBs je ein Factsheet mit bereichsspezifischen Zielen und $\gamma$-Zielen
    \item \textbf{GL-Workshop:} Validierung des Gesamtmodells -- inkl. Simulation verschiedener Implementierungsszenarien
    \item \textbf{Strategie-Handbuch:} Vollständige Dokumentation (20--30 Seiten) für die laufende Nutzung
\end{itemize}

\textbf{Deliverables:} $\gamma$-Dashboard, 5 Bereichs-Factsheets, Strategie-Handbuch

\subsection{Meilensteine und Zeitplan}

\begin{center}
\begin{tabular}{p{2cm}p{3cm}p{5.5cm}p{2.5cm}}
\toprule
\rowcolor{fadarkblue}
\fahead{Woche} & \fahead{Phase} & \fahead{Meilenstein} & \fahead{Typ} \\
\midrule
\rowcolor{fawhite}
1 & Konsolidierung & Kick-off mit CEO + Gap-Präsentation & -- \\
\rowcolor{falightgray}
2--3 & Konsolidierung & Bereichsleiter-Gespräche (5$\times$) & -- \\
\rowcolor{fawhite}
4 & Konsolidierung & Strategie-Pyramide (GL-Freigabe) & SOFT \\
\rowcolor{falightgray}
8 & Modellentwicklung & Modell-Präsentation GL & HARD \\
\rowcolor{fawhite}
9--10 & Implementierung & $\gamma$-Dashboard + Factsheets & SOFT \\
\rowcolor{falightgray}
11 & Implementierung & GL-Workshop Validierung & SOFT \\
\rowcolor{fawhite}
12 & Implementierung & VR-Präsentation + Handbuch & HARD \\
\bottomrule
\end{tabular}
\end{center}

\subsection{Was wir von der SOB brauchen}

\begin{itemize}
    \item \textbf{CEO:} 4--6 Stunden pro Monat (Steering)
    \item \textbf{GL-Mitglieder:} 2 Workshops (je ca. 2h)
    \item \textbf{Bereichsleiter (5):} Je 1 Gespräch (90 Min.) + 1 Factsheet-Review (60 Min.)
    \item \textbf{Dokumentenzugang:} Bereichsstrategien, Budgets, KPI-Reports (falls vorhanden)
\end{itemize}

\newpage
% =============================================================================
\section{Unser Leistungsangebot}
% =============================================================================

\subsection{Investition}

\begin{center}
\fcolorbox{fadarkblue}{falightgray}{%
\parbox{0.85\textwidth}{%
\centering
\vspace{0.5cm}
\textbf{SOB002 -- Integriertes Strategiemodell}\\[0.3cm]
\textbf{Dauer:} 3 Monate (Februar -- April 2026)\\[0.3cm]
\textbf{\Large CHF 30'000}\\[0.3cm]
\small{Fixpreis, pauschal}\\[0.2cm]
\small{exkl. 8.1\% MwSt.}\\[0.2cm]
\small{Reisespesen nach Aufwand}
\vspace{0.5cm}
}%
}
\end{center}

\subsection{Im Preis inbegriffen}

\begin{center}
\begin{tabular}{p{0.5cm}p{5cm}p{3cm}p{3cm}}
\toprule
\rowcolor{fadarkblue}
\fahead{\#} & \fahead{Deliverable} & \fahead{Format} & \fahead{Termin} \\
\midrule
\rowcolor{fawhite}
1 & Strategie-Pyramide (vollständig) & PDF (1 Seite) & Woche 4 \\
\rowcolor{falightgray}
2 & SOB-Strategiemodell & BEATRIX + YAML & Woche 8 \\
\rowcolor{fawhite}
3 & Dimensions-Katalog & YAML + Excel & Woche 8 \\
\rowcolor{falightgray}
4 & $\gamma$-Dashboard & BEATRIX & Woche 10 \\
\rowcolor{fawhite}
5 & Bereichs-Factsheets (5$\times$) & PDF (je 2 S.) & Woche 11 \\
\rowcolor{falightgray}
6 & Strategie-Handbuch & PDF (20--30 S.) & Woche 12 \\
\bottomrule
\end{tabular}
\end{center}

\vspace{0.5cm}

Zusätzlich inbegriffen:
\begin{itemize}
    \item GL-Workshop zur Validierung des Gesamtmodells
    \item 5 Bereichsleiter-Gespräche (je 60--90 Minuten)
    \item Laufende Abstimmung mit CEO (Steering, 4--6h/Monat)
    \item VR-Präsentation zum Projektabschluss
    \item \textbf{Vorarbeiten (bereits erbracht):} 11-Dokument-Analyse, Gap-Analyse (47 Gaps), $\gamma$-Matrix
\end{itemize}

\subsection{Nicht inbegriffen}

\begin{itemize}
    \item Detaillierte Massnahmenplanung pro Bereich (mögliches Folge-Projekt SOB003)
    \item Change Management / Kommunikationskonzept
    \item IT-Implementierung ausserhalb BEATRIX
    \item Schulungen / Trainings für Führungskräfte
    \item Quantitative Marktforschung / Kundenbefragung
    \item Lohn-/Anreizsystem-Redesign
\end{itemize}

\subsection{Zahlungskonditionen}

\begin{itemize}
    \item 40\% bei Projektstart (CHF 12'000)
    \item 30\% nach Phase 2 / Modell-Präsentation GL (CHF 9'000)
    \item 30\% nach Projektabschluss / VR-Präsentation (CHF 9'000)
\end{itemize}

Zahlungsfrist: 30 Tage netto

\subsection{Gültigkeit und AGB}

Diese Offerte ist gültig bis \textbf{31. März 2026}. Es gelten die Allgemeinen Geschäftsbedingungen der FehrAdvice \& Partners AG in der aktuellen Fassung.

\newpage
% =============================================================================
\section*{Annahme}
% =============================================================================

Mit der Unterzeichnung dieser Offerte erteilt die Schweizerische Südostbahn AG den Auftrag für das Projekt SOB002 <<SOB Strategiemodell>> gemäss den oben aufgeführten Bedingungen.

\vspace{2cm}

\begin{center}
\begin{tabular}{p{7cm}p{7cm}}
\textbf{Auftraggeber} & \textbf{Auftragnehmer} \\[0.5cm]
Schweizerische Südostbahn AG & FehrAdvice \& Partners AG \\[2cm]
\rule{6cm}{0.4pt} & \rule{6cm}{0.4pt} \\
Armin Weber, CEO & Gerhard Fehr, CEO \\[1cm]
\rule{6cm}{0.4pt} & \rule{6cm}{0.4pt} \\
Ort, Datum & Ort, Datum \\
\end{tabular}
\end{center}

\vspace{2cm}

\begin{center}
\small{
FehrAdvice \& Partners AG | Klausstrasse 4 | 8008 Zürich\\
Tel. +41 44 256 79 00 | info@fehradvice.com | www.fehradvice.com
}
\end{center}

\end{document}
